\documentclass[10pt,a4paper]{amsart}
\usepackage[margin=19mm]{geometry}
\usepackage{amsmath,amssymb,mathtools}
\usepackage[scr=rsfs,cal=euler]{mathalfa}
\usepackage{xfrac}
\usepackage[unicode,hidelinks,bookmarks=false]{hyperref}
\newcommand{\ie}{i.e.\ }
\newcommand{\cf}{cf.\ }
\newcommand{\resp}{respectively\ }
\newcommand{\Subsection}{\subsection}
\providecommand{\nobreakdash}{\nobreak\hbox{-}}
\setlength{\emergencystretch}{2em}
\multlinegap0pt
\usepackage[shortlabels]{enumitem}
\usepackage{amssymb}
\usepackage{mathtools}
\usepackage{tensor}
\newtheorem{definition}{Definition}
\newtheorem{lemma}{Lemma}
\newtheorem{claim}{Claim}
\DeclareMathOperator{\Fn}{Fn}
\newcommand{\I}{\mathcal{I}}
\renewcommand{\a}{\alpha}
\DeclareMathOperator{\dom}{dom}
\renewcommand{\k}{\kappa}
\newcommand{\markdef}[1]{\textbf{#1}}
\newcommand{\pre}[2]{\tensor[^{#1}]{#2}{}}
\newcommand{\seq}[2]{\langle #1 \mid #2 \rangle}
\title{On Borel sets in ideal topologies}
\author{Miguel Moreno, Beatrice Pitton}
\date{Source: arXiv:2512.21140v3 (CC BY 4.0)}
\begin{document}
\maketitle

\noindent\emph{Source and licence.} On Borel sets in ideal topologies. Version arXiv:2512.21140v3 (CC BY 4.0), distributed by arXiv under the Creative Commons Attribution 4.0 International licence, http://creativecommons.org/licenses/by/4.0/, as recorded in the arXiv licence metadata for this exact version and shown on the abstract page https://arxiv.org/abs/2512.21140v3. Full licence notice: \texttt{licenses/arxiv-2512.21140-CC-BY-4.0.txt}.\par\smallskip

\noindent\textit{Editorial scope.} Counted unit: the article's lemma approx\_lemma (source0.tex lines 1236--1242). The article's complete original proof of it is delivered with it (source0.tex lines 1244--1301, one \texttt{proof} environment of 58 lines). No mathematics is written, repaired or re-derived: the statement and the proof are the article's own lines, in the article's own notation, and no retained line is substituted or re-flowed.  Retained uncounted as the source-local context the proof needs: lines 1221--1229.  Nothing was omitted from any retained span, so the package carries no omission record.  No retained line contains a citation, so the wrapper carries no bibliography and no bibliography entry.  No retained line contains a figure environment, an image-inclusion command or drawing code, so no image is delivered and the package declares no asset dependency.  Editorial formatting only, in addition: an AMS \texttt{amsart} preamble that restates the article's own declarations and macros used by the retained lines, namely the environments \texttt{definition}, \texttt{lemma}, \texttt{claim} on separate counters, together with the standard packages those lines need.  The retained environments print in this document as Definition 1, Lemma 1, Claim 1 in order of appearance, whereas the article prints the same items with the numbering of its own full document.  The complete native source of the article as distributed in the arXiv e-print for this exact version is bundled unchanged as \texttt{sources/191619/source0.tex}.
\begin{definition}\label{A-approximationk}
    Let \(\seq{B_\alpha}{\alpha<\kappa}\) be an approximation sequence of $\mathcal{I}$, and let  $\mathcal{H}$ be a filter over $\kappa$ with basis $\mathcal{B}$. Let $A\in \mathcal{H}^+$ be an unbounded set and let $E$ be an equivalence relation on $\pre{\k}{\kappa}$.
    We say that $E$ has an \markdef{$A$-approximation over $\seq{B_\alpha}{\alpha<\kappa}$} if there is  a sequence $\langle E_\alpha\mid \alpha<\kappa\rangle$ with $E_\alpha\subseteq \pre{B_{s(\alpha)}}{\k}\times \pre{B_{s(\alpha)}}{\k}$ such that the following conditions hold:
    \begin{enumerate}[label={\upshape (\arabic*)}, leftmargin=2pc]
        \item For every $ \alpha<\kappa$, $E_\alpha$ is an equivalence relation.
        \item\label{A-approximationk_2} For every $x,y\in \pre{\k}{\kappa}$, if $x E y$ then there exists $D\in \mathcal{B}$ such that for every $\alpha\in D$, $x\restriction B_{s(\alpha)} \ E_\alpha\ y\restriction B_{s(\alpha)}$.
        \item\label{A-approximationk_3} For every $x,y\in \pre{\k}{\kappa}$, if $\neg (x E y)$ then there exists $A'\in \mathcal{H}^+$, $A'\subseteq A$, such that for every $\alpha\in A'$, $\neg(x\restriction B_{s(\alpha)} \ E_\alpha\ y\restriction B_{s(\alpha)})$.
    \end{enumerate}
\end{definition}

\begin{lemma}[Approximation Lemma]\label{approx_lemma}
    Let $\mathcal{H}$ be a filter over $\kappa$ with basis $\mathcal{B}$. Let $A\in \mathcal{H}^+$ be an unbounded set and let $E$ be an equivalence relation on $\pre{\k}{\kappa}$. Then, the following are equivalent:
    \begin{enumerate}[label={\upshape (\arabic*)}, leftmargin=2pc]
        \item\label{approx1} $E$ has a precise $A$-approximation;
        \item\label{approx2} $E\hookrightarrow_{R(A)}\sim^\kappa_\mathcal{F}$, where $\mathcal{F}=\mathcal{H}\restriction A$.
\end{enumerate}
\end{lemma}

\begin{proof}
\ref{approx1}$\Rightarrow$ \ref{approx2}.
Let $\seq{B_\alpha}{\alpha<\kappa}$ be an approximation sequence of $\mathcal{I}$ such that there is $\langle E_\alpha\mid \alpha<\kappa\rangle$ a precise $A$-approximation over $\seq{B_\alpha}{\alpha<\kappa}$. For every $\alpha<\kappa$, let $\langle x^\alpha_i\mid 0<i<\kappa\rangle$ be an enumeration of the equivalence classe $E_\alpha$. 
    Let us define $F:\pre{\k}{\kappa}\rightarrow\pre{\k}{\kappa}$ as follows:
        $$F(x)(\alpha)=\begin{cases} i &\mbox{if }\alpha\in A \mbox{ and } x\restriction B_{s(\alpha)}\in x_i^{\alpha},\\
    0 & \mbox{otherwise. }\end{cases}$$  
    Let us show that $x\ E\ y$ if and only if $F(x)\ \sim^\kappa_\mathcal{F}\ F(y)$.
    \begin{claim}
        $x\ E\ y$ implies $F(x)\ \sim^\kappa_\mathcal{F}\ F(y)$.
    \end{claim}
    \begin{proof}[Proof of the Claim.]
        Suppose $x,y\in \pre{\k}{\kappa}$ are such that $x\ E\ y$. Since $\langle E_\alpha\mid \alpha<\kappa\rangle$ is an $A$-approximation over $\seq{B_\alpha}{\alpha<\kappa}$, by Definition~\ref{A-approximationk}\ref{A-approximationk_2} there exists $D\in\mathcal{B}$ such that for every $\alpha\in D$, $$x\restriction B_{s(\alpha)}\ E_\alpha\ y\restriction B_{s(\alpha)}.$$ So, for all $\alpha\in D\cap A$, $F(x)(\alpha)=F(y)(\alpha)$. Since $D\cap A\in \mathcal{H}\restriction A=\mathcal{F}$, we get $F(x)\ \sim^\kappa_\mathcal{F}\ F(y)$.
    \end{proof}
    \begin{claim}
        $\neg(x\ E\ y)$ implies $\neg(F(x)\ \sim^\kappa_\mathcal{F}\ F(y))$.
    \end{claim}
    \begin{proof}[Proof of the Claim.]
        Suppose $x,y\in \pre{\k}{\kappa}$ are such that $\neg(x\ E\ y)$. Since $\langle E_\alpha\mid \alpha<\kappa\rangle$ is an $A$-approximation over $\seq{B_\alpha}{\alpha<\kappa}$, by Definition~\ref{A-approximationk}\ref{A-approximationk_3}, there is $A'\subseteq A$ in $\mathcal{H}^+$,  such that for all $\alpha\in A'$, $$\neg(x\restriction B_{s(\alpha)}\ E_\alpha\ y\restriction B_{s(\alpha)}).$$ So, for all $\alpha\in A'$, $F(x)(\alpha)\neq F(y)(\alpha)$. Since $A'\in\mathcal{H}^+$, for all $X\in \mathcal{H}$, $X\cap A'\neq\emptyset$. Thus, for every $X\in \mathcal{H}$, $(X\cap A)\cap A'\neq\emptyset$. So $A'\in (\mathcal{H}\restriction A)^+=\mathcal{F}^+$, and
        we conclude that $\neg(F(x)\ \sim^\kappa_\mathcal{F}\ F(y))$.
    \end{proof}
    \begin{claim}
        $F$ is $A$-recursive over $\seq{B_\alpha}{\alpha<\kappa}$.
    \end{claim}
    \begin{proof}[Proof of the Claim.]
    Define $H:\Fn_{\I}\rightarrow \Fn_{\I}$ as follows:
        $$H(x\restriction D)=F(x)\restriction s^{-1}[\alpha+1]$$%\begin{cases} F(x)\restriction B_{\alpha'}\cup (A\cap \alpha') &\mbox{if }D=B_\alpha, \alpha\in A, \mbox{ and } \alpha'=min(A\backslash (\alpha+1))$$; \\
    %\bar0_\alpha & \mbox{otherwise. }\end{cases}$$ 
   % Clearly, if $\beta<\alpha$, then $H(x\restriction B_\beta)\subseteq H(x\restriction B_\alpha)$. 
    where $\alpha$ is the smallest element such that $D\subseteq B_\alpha$.
    Notice that if $D=B_\alpha$, then $H(x\restriction D)=H(x\restriction B_\alpha)= F(x)\restriction s^{-1}[\alpha+1]$. Also, for all $\beta<\kappa$, $\beta\in s^{-1}[s(\beta)+1]$ and $H(x\restriction B_{s(\beta)})(\beta)=F(x)(\beta)$. Finally $F$ is 0 outside $A$. 

   Thus we are only missing to show that $H$ is well defined, it is sufficient to check the case $\dom(x)=B_\alpha$. Let $x,y\in \kappa^{\kappa}$ be such that $x\restriction B_\alpha=y\restriction B_\alpha$. 
    Clearly, for every $\beta\in s^{-1}[\alpha+1]\backslash A$, $F(x)(\beta)=0=F(y)(\beta)$, therefore $F(x)\restriction s^{-1}[\alpha+1]\ (\beta)=0=F(y)\restriction s^{-1}[\alpha+1]\ (\beta)$ for all $\beta\in s^{-1}[\alpha+1]\backslash A$.
    On the other hand, let $\beta\in s^{-1}[\alpha+1]\cap A$. From the definition of $F$, $F(x)(\beta)=i$ and $F(y)(\beta)=j$, where $x\restriction B_{s(\beta)}\in x_i^\beta$ and $y\restriction B_{s(\beta)}\in x_j^\beta$. Since $\beta\in s^{-1}[\alpha+1]\cap A$, $s(\beta)\leq \alpha$. Thus $x\restriction B_{s(\beta)}=y\restriction B_{s(\beta)}$, $x_i^\beta=x_j^\beta$, and $i=j$. 
    So $F(x) (\beta)=F(y) (\beta)$ for all $\beta\in s^{-1}[\alpha+1]\cap A$. We conclude that $F(x)\restriction s^{-1}[\alpha+1] =F(y)\restriction $, hence $H(x\restriction B_\alpha)=H(y\restriction B_\alpha)$ and $H$ is well defined.

   Finally, $F$ is $A$-recursive by definition of \(H\). 
    \end{proof}
    %Since for all $x$, $F(x)$ is constant to 0 outside $C\cap A$, $F$ is $S$-recursive for all $S\supset A$. In particular, $F$ is $\kappa$-recursive, thus Lipschitz.
    
\ref{approx2}$\Rightarrow$ \ref{approx1}.
Let $\seq{B_\alpha}{\alpha<\kappa}$ be an approximation sequence of $\mathcal{I}$ and $f:\pre{\k}{\kappa}\rightarrow\pre{\k}{\kappa}$ be an $A$-recursive function over $\seq{B_\alpha}{\alpha<\kappa}$ witnessing $E\ \hookrightarrow_{R(A)}\ \sim^\kappa_\mathcal{F}$. 
    For every $\alpha<\kappa$, define $E_\alpha\subseteq \pre{B_{s(\alpha)}}{\k}\times\pre{B_{s(\alpha)}}{\k}$ as follows:
    \begin{itemize}
        \item if $\alpha\notin A$, let $E_\alpha=\pre{B_{s(\alpha)}}{\k}\times\pre{B_{s(\alpha)}}{\k}$ (notice that $\alpha\notin A$ implies $f(x)(\alpha)=0$);
        \item if $\alpha\in A$, let \(E_\a\) be such that $x\ E_\alpha\ y$ if and only if $f(x^0)(\alpha)=f(y^0)(\alpha)$.
   \end{itemize}
%    \begin{claim}
     We now show that $\langle E_\alpha\mid \alpha<\kappa\rangle$ is a precise $A$-approximation over $\seq{B_\alpha}{\alpha<\kappa}$.
%    \end{claim}
It is clear that $E_\alpha$ is an equivalence relation for every $\alpha<\k$, with at most $\k$ many classes. To see that $\langle E_\alpha\mid \alpha<\kappa\rangle$ satisfies Definition~\ref{A-approximationk}\ref{A-approximationk_2}, let $x,y\in\pre{\k}{\kappa}$  be such that $x Ey$. Since $f$ is an $A$-reduction from $E$ to $\sim^\kappa_\mathcal{F}$, $f(x)\sim^\kappa_\mathcal{F} f(y)$. Thus there is $W\in \mathcal{F}$ such that $W\subseteq \{\alpha<\kappa\mid f(x)(\alpha)= f(y)(\alpha)\}$. 
Since $\mathcal{F}=\mathcal{H}\restriction A$, there is $D\in \mathcal{B}$ such that $D\cap A\subseteq\{\alpha<\kappa\mid f(x)(\alpha)= f(y)(\alpha)\}$.
On the other hand $f$ is $A$-recursive over $\seq{B_\alpha}{\alpha<\kappa}$, there is $H:\Fn_{\I}\rightarrow \Fn_{\I}$ such that for all $\alpha\in \kappa$, $$f((x\restriction B_{s(\alpha)})^0)(\alpha)=H((x\restriction B_{s(\alpha)})^0\restriction B_{s(\alpha)})(\alpha)=H(x\restriction B_{s(\alpha)})(\alpha)=f(x)(\alpha)$$ and $$f(y)(\alpha)=f((y\restriction B_{s(\alpha)})^0)(\alpha).$$ 
%From the definition of $s$, we know that $s(\alpha)\leq \alpha$, thus $$f(x)(\alpha)=f((x\restriction B_{s(\alpha)})^0)(\alpha) \text{ and } f(y)(\alpha)=f((y\restriction B_{s(\alpha)})^0)(\alpha).$$ 
Since for all $\alpha\in D\cap A$, $f(x)(\alpha)=f(y)(\alpha)$, $f((x\restriction B_{s(\alpha)})^0)(\alpha)=f((y\restriction B_{s(\alpha)})^0)(\alpha)$ holds for all $\alpha\in D\cap A$. We conclude that for all $\alpha\in D\cap A$, $x\restriction B_{s(\alpha)}\ E_\alpha\ y\restriction B_{s(\alpha)}$. 

We conclude by showing that $\langle E_\alpha\mid \alpha<\kappa\rangle$ satisfies Definition~\ref{A-approximationk}\ref{A-approximationk_3}. Let $x,y\in\pre{\k}{\kappa}$ be such that $\neg(x\ E\ y)$. There is $Y\in \mathcal{F}^+$ such that for all $\alpha\in Y$, $f(x)(\alpha)\neq f(y)(\alpha)$. Since $\mathcal{F}=\mathcal{H}\restriction A$, $A\cap Y\in \mathcal{F}^+$. Let us denote $A\cap Y$ by $A'$. Therefore, $A'\subseteq A$ is such that for all $\alpha\in A'$, $f(x)(\alpha)\neq f(y)(\alpha)$. Using a similar argument as above, we conclude that $f((x\restriction B_{s(\alpha)})^0)(\alpha)\neq f((y\restriction B_{s(\alpha)})^0)(\alpha)$ holds for all $\alpha\in A'$. We conclude that for all $\alpha\in A'$, $\neg (x\restriction B_{s(\alpha)}\ E_\alpha\ y\restriction B_{s(\alpha)})$.
\end{proof}

\end{document}
